\chapter{RL Meets the Real World}
\label{ch:rl-real-world}

Reinforcement learning does not exist in a vacuum.
It can be trained from scratch, layered on top of supervised-learning base models, woven into diffusion and flow-matching generators, and shuttled between simulation and reality.
This chapter maps the landscape of how RL connects to other paradigms---and where each connection matters most.

%------------------------------------------------------------------
\section{Pure RL from Scratch}
\label{sec:pure-rl}

The classic setting: no pretrained model, random initialisation, learning purely from environment interaction.

\subsection{Value-Based Methods}

The model learns a scalar estimate of long-term return for every state--action pair:
\[
  Q(s,a) \;\approx\; \mathbb{E}\!\Bigl[\sum_{k=0}^{\infty}\gamma^k r_{t+k} \;\Big|\; s_t=s,\;a_t=a\Bigr].
\]
Decisions follow $a = \arg\max_a Q(s,a)$.
Representatives: Q-learning, DQN, Double DQN, Dueling DQN.

Essence: \emph{learn valuations first, then extract behaviour from valuations}.

\subsection{Policy-Based Methods}

The model directly learns a preference distribution over actions:
\[
  \pi_\theta(a|s) \;\to\; \text{behavioural preference distribution}.
\]
Decisions follow $a \sim \pi_\theta(\cdot|s)$.
Representatives: REINFORCE, PPO, SAC.

Essence: \emph{directly learn the preference distribution}, skipping the intermediate value landscape.

\subsection{Actor-Critic: Combining Both}

Most modern algorithms are Actor-Critic:
\begin{itemize}
  \item \textbf{Critic}: learns value ($V$ or $Q$), provides a preference signal (advantage).
  \item \textbf{Actor}: learns the behaviour distribution, guided by the Critic.
\end{itemize}
Representatives: A2C, DDPG, TD3, SAC, PPO (with a $V$-critic).

%------------------------------------------------------------------
\section{Architecture: Training vs.\ Inference Modules}
\label{sec:architecture}

A recurring pattern across all RL architectures: \emph{training requires more modules than inference}.

\begin{keyidea}[Critic = Scaffolding]
The Critic (or Value Head) exists only during training.
It provides gradient signals---advantage estimates or $Q$-gradients---that guide the Actor's updates.
At deployment, the Critic is discarded; only the Actor/Policy network is shipped to production.
\end{keyidea}

\begin{table}[htbp]
\centering
\caption{What is used at training vs.\ inference.}
\begin{tabular}{@{}llll@{}}
\toprule
\textbf{Architecture} & \textbf{Training} & \textbf{Inference} & \textbf{Discarded} \\
\midrule
DQN & Backbone + Q Head + Target Net & Backbone + Q Head & Target Net \\
PPO (shared) & Backbone + Policy Head + Value Head & Backbone + Policy Head & Value Head \\
SAC / TD3 & Actor Net + Critic Net(s) & Actor Net & Entire Critic \\
\bottomrule
\end{tabular}
\end{table}

%------------------------------------------------------------------
\section{Reward Sources: Prepared Before Training}
\label{sec:reward-sources}

The RL loop requires reward signals, but reward is \emph{not learned during RL training}---it must be prepared beforehand.

\begin{enumerate}
  \item \textbf{Environment-provided reward.}
    Game scores, distance, collision penalties, energy consumption.
    The environment directly returns a scalar; no extra preparation needed.

  \item \textbf{Rule-based reward.}
    Math-answer correctness, code compilation pass rate, format compliance.
    A hand-written function $R(\text{output})\to\mathbb{R}$.
    Common with GRPO.

  \item \textbf{Learned Reward Model (RM).}
    Collect human preference pairs (chosen, rejected), train an SL model
    $\text{RM}(x,y)\to\mathbb{R}$ to approximate human preferences, then freeze it.
    Standard approach for PPO-RLHF (InstructGPT).
\end{enumerate}

\begin{keyidea}[RM Is a Precondition --- Not Part of RL]
The Reward Model is trained via supervised learning on preference data and frozen as a ``scorer'' during RL training.
It is logically \emph{upstream} of RL, not inside it.
\end{keyidea}

%------------------------------------------------------------------
\section{Value and Policy: Coupled in Sampling, Decoupled in Gradients}
\label{sec:coupled-decoupled}

In Actor-Critic training, the Critic and Actor appear in the same loop, but they are \textbf{two independent optimisation processes}:

\[
\begin{aligned}
  \mathcal{L}_{\text{critic}} &= \bigl(V_\phi(s) - V_{\text{target}}\bigr)^2,
  \qquad &\text{updates only } \phi; \\
  \mathcal{L}_{\text{actor}}  &= -\log\pi_\theta(a|s)\cdot A(s,a),
  \qquad &\text{updates only } \theta.
\end{aligned}
\]

\textbf{Sampling level: coupled.}
Policy determines what trajectories are collected; those trajectories are the data for both Critic and Actor.
Change the policy, and the data distribution changes for everyone.

\textbf{Gradient level: decoupled.}
When computing $\nabla_\theta \mathcal{L}_{\text{actor}}$, the advantage $A(s,a)$ is treated as a \emph{constant}---a \texttt{stop\_gradient} wall separates the two.
The Critic's parameters receive no gradient from the Actor's loss, and vice versa.

\[
  \nabla_\theta J = \mathbb{E}\!\bigl[\nabla_\theta \log\pi_\theta(a|s)\cdot \underbrace{A(s,a)}_{\text{detached}}\bigr]
\]

The one exception is DDPG/TD3's ``max-$Q$ loss'' $\mathcal{L}_{\text{actor}} = -Q(s,\mu_\theta(s))$, where the gradient flows \emph{through} $Q$ to the action and then to the actor---but $Q$'s own parameters are still frozen during the actor update.

This alternating-iteration structure---Critic chases the policy's changing data, Actor chases the Critic's changing landscape---is the root cause of Actor-Critic instability.

%------------------------------------------------------------------
\section{RL Fine-Tuning of SL Base Models}
\label{sec:rl-finetune}

The dominant paradigm for LLM alignment: pretrain with SL, then fine-tune with RL.

\[
  \text{Pretraining (SL)} \;\to\; \text{SFT (SL)} \;\to\; \text{Preference Fine-Tuning (RL)}
\]

\subsection{PPO Approach (InstructGPT Route)}

Requires \textbf{four models} in memory simultaneously:
\begin{enumerate}
  \item \textbf{Policy Model} --- the LLM being trained.
  \item \textbf{Reference Model} --- frozen SFT copy, for KL constraint.
  \item \textbf{Reward Model} --- trained from preferences, frozen as scorer.
  \item \textbf{Critic / Value Head} --- estimates $V(s)$ for advantage computation; typically initialised from the RM.
\end{enumerate}

Training loop: Policy generates $\to$ RM scores $\to$ Critic estimates baseline $\to$ compute advantage $\to$ PPO clipped-ratio update $\to$ KL penalty keeps Policy near Reference.

\subsection{GRPO Approach (DeepSeek Route)}

Eliminates the Critic entirely.
Requires only \textbf{two--three models}: Policy, Reference, and a reward source (RM or rule-based verifier).

Core mechanism:
\begin{enumerate}
  \item For each prompt, Policy generates a \emph{group} of $G$ responses.
  \item Each response is scored.
  \item Within-group relative ranking serves as baseline (no Critic needed).
  \item Policy gradient updates the Policy.
\end{enumerate}

DeepSeek-R1 uses full fine-tuning with GRPO, training all parameters.

\subsection{DPO Approach (Stanford Route)}

Eliminates even the Reward Model.
Requires only \textbf{two models}: Policy and Reference.
Trains directly on preference pairs $(\text{chosen}, \text{rejected})$ with a classification-style loss:
\[
  \mathcal{L}_{\text{DPO}} = -\log\sigma\!\Bigl(\beta\Bigl[
    \log\frac{\pi_\theta(\text{chosen})}{\pi_{\text{ref}}(\text{chosen})}
    - \log\frac{\pi_\theta(\text{rejected})}{\pi_{\text{ref}}(\text{rejected})}
  \Bigr]\Bigr).
\]

\subsection{Comparison}

\begin{table}[htbp]
\centering
\caption{LLM RL fine-tuning approaches.}
\begin{tabular}{@{}lcccc@{}}
\toprule
\textbf{Approach} & \textbf{Models} & \textbf{Bellman eq.} & \textbf{Reward source} & \textbf{Complexity} \\
\midrule
PPO & 4 & Used (Critic) & RM & High \\
GRPO & 2--3 & Not used & RM / rules & Medium \\
DPO & 2 & Not used & Pref.\ pairs & Low \\
\bottomrule
\end{tabular}
\end{table}

\subsection{Parameter-Efficient Updates: LoRA and QLoRA}

Regardless of the RL method, one can choose \emph{how} to update the Policy:

\begin{itemize}
  \item \textbf{Full fine-tuning}: update all parameters. Maximum expressiveness; huge memory.
  \item \textbf{LoRA}: freeze the base, train low-rank matrices $\Delta W = AB$. The Reference Model is free---a forward pass without the adapter \emph{is} the reference.
  \item \textbf{QLoRA}: 4-bit quantised base + LoRA. Enables RLHF on consumer GPUs.
\end{itemize}

Typical choices: PPO with full fine-tuning (InstructGPT) or LoRA; GRPO with full fine-tuning (DeepSeek-R1); DPO + LoRA as the most memory-efficient alignment path.

%------------------------------------------------------------------
\section{RL with Diffusion and Flow Matching Models}
\label{sec:diffusion-rl}

When a diffusion or flow-matching model is used as a \emph{policy}, the multi-step denoising/flow process can be cast as an MDP and fine-tuned with RL.

\subsection{Denoising Chain as MDP}

\begin{keyidea}[DDPO: Denoising Diffusion Policy Optimisation]
Model the denoising chain as an episodic MDP:
\emph{state} = current noisy sample $x_t$;
\emph{action} = one denoising step's output;
\emph{episode} = complete chain $x_T \to x_0$;
\emph{reward} = given only on the final $x_0$.
Apply REINFORCE or PPO to estimate the policy gradient for each step without backpropagating through the entire chain.
\end{keyidea}

DDPO (Black et al., 2023) successfully fine-tuned text-to-image diffusion models with aesthetic scores and human preferences.

\textbf{DPPO} (Ren et al., 2025) adapts PPO specifically for diffusion policies in continuous control, addressing cross-step clipping and high-dimensional action spaces.

\textbf{DRaFT} (Clark et al., 2023) backpropagates reward gradients directly through a truncated denoising chain---more sample-efficient but requires a differentiable RM.

\subsection{Flow Matching and RL}

Flow Matching replaces stochastic SDEs with deterministic ODEs, enabling 5--10$\times$ faster inference.
For RL, each discretised ODE step is treated as an action in the MDP, mirroring the DDPO formulation.

\subsection{Key Challenges}

\begin{table}[htbp]
\centering
\caption{Challenges in RL + diffusion/flow.}
\begin{tabular}{@{}ll@{}}
\toprule
\textbf{Challenge} & \textbf{Common solution} \\
\midrule
Cross-step credit assignment & Episodic MDP formulation \\
High REINFORCE variance & Fewer steps / baselines / truncated backprop \\
Reward differentiability & Policy gradient bypasses this \\
KL regularisation & Penalty similar to RLHF \\
Action-chunk consistency & Overlapping window design \\
\bottomrule
\end{tabular}
\end{table}

%------------------------------------------------------------------
\section{Agent RL vs.\ One-Round RL}
\label{sec:agent-rl}

\subsection{Traditional RL: Single-Round Sequential Interaction}

A continuous time-series rollout:
$s_0 \to a_0 \to r_0 \to s_1 \to a_1 \to r_1 \to \cdots \to s_T$.
Each action directly produces the next state and reward.
Typical scenarios: autonomous driving, robot control, Atari games.

\subsection{Agentic RL: Multi-Round Conversational Interaction}

An episode is a multi-round dialogue:
\begin{enumerate}
  \item Agent observes $\to$ thinks $\to$ calls tools / outputs text $\to$ environment feedback.
  \item Agent observes feedback $\to$ thinks $\to$ acts again $\to$ feedback.
  \item[$\vdots$]
  \item Agent produces final output $\to$ episode ends.
\end{enumerate}
Typical scenarios: coding agents, dialogue systems, tool-using agents.

\subsection{The Fundamental Connection}

\begin{keyidea}[Autonomous Driving Is Agentic RL]
In autonomous driving, each action chunk (a control sequence spanning $\sim$100\,ms--2\,s) is analogous to one ``round'' in an agentic AI system.
After each chunk, the agent observes the new environment state and decides again.
Both are closed-loop MDPs; the differences are only in action granularity, environment nature, and interaction frequency.
\end{keyidea}

\begin{table}[htbp]
\centering
\caption{Traditional RL vs.\ Agentic RL.}
\begin{tabular}{@{}lll@{}}
\toprule
\textbf{Dimension} & \textbf{Autonomous Driving} & \textbf{Agentic AI} \\
\midrule
One ``round'' & Action chunk (100\,ms--2\,s) & Agent output (tool call / text) \\
Env.\ feedback & New sensor observations & Tool results / user replies \\
Episode & One drive (min--hrs) & One task (code fix / Q\&A) \\
State space & Continuous, high-dim & Discrete, variable-length \\
Action space & Continuous & Discrete / mixed \\
\bottomrule
\end{tabular}
\end{table}

\subsection{Training Differences}

\textbf{Traditional RL} (driving / robotics): continuous-action policies (Gaussian / Diffusion / Flow), Actor-Critic (SAC / TD3 / PPO), high-frequency simulator interaction, engineered rewards.

\textbf{Agentic RL} (LLM agents): autoregressive token policy, PPO / GRPO / DPO, multi-round agent--environment interaction, task-completion / preference / verifier rewards, far higher compute per ``action.''

%------------------------------------------------------------------
\section{Sim2Real / Real2Sim Loop}
\label{sec:sim2real}

RL needs massive interaction, but real-world interaction is expensive, dangerous, and slow.
The Sim2Real and Real2Sim loop addresses this.

\subsection{Sim2Real: From Simulation to Reality}

The core challenge is the \textbf{reality gap}---a policy trained in simulation often fails in the real world because physics, rendering, and dynamics differ.

\textbf{Domain Randomisation.}
Randomise simulator parameters (friction, mass, lighting, textures, sensor noise) during training so the real world is ``just another draw'' from the training distribution.
OpenAI's Automatic Domain Randomisation (ADR, 2019) trained a robotic hand to solve a Rubik's cube by dynamically expanding randomisation across thousands of parameters.

\textbf{System Identification.}
Measure real-world physics parameters, calibrate the simulator, and optionally randomise \emph{around} calibrated values rather than uniformly.

\textbf{Domain Adaptation.}
Use CycleGAN or adversarial training to align visual distributions between sim and real without paired data.

\textbf{Progressive Transfer.}
Train in simulation first, then fine-tune in the real environment with a small amount of real data.

\subsection{Real2Sim: From Reality to Simulation}

The reverse direction: use real-world data to build or improve simulators.

\begin{itemize}
  \item \textbf{Neural Scene Reconstruction}: NeRF / 3D Gaussian Splatting reconstruct photo-realistic, editable scenes from real sensor data.
  \item \textbf{Learned World Models}: generative models trained on real data serve as implicit simulators (GAIA-1, DayDreamer).
  \item \textbf{Digital Twins}: high-fidelity digital replicas continuously updated with sensor data.
  \item \textbf{Sim Parameter Calibration}: optimise simulator parameters (friction, drag, latency) using real trajectory data via Bayesian optimisation or differentiable simulation.
\end{itemize}

\subsection{The Continuous Loop}

In autonomous driving, the loop runs continuously:
\[
  \text{Real data} \xrightarrow{\text{Real2Sim}} \text{Better sim}
  \xrightarrow{\text{Train}} \text{Better policy}
  \xrightarrow{\text{Sim2Real}} \text{Deploy}
  \xrightarrow{\text{Collect}} \text{New real data}
  \to \cdots
\]
Each iteration reduces the reality gap and improves real-world performance.

%------------------------------------------------------------------
\section{Summary}

\begin{table}[htbp]
\centering
\caption{RL's role across different scenarios.}
\begin{tabular}{@{}lll@{}}
\toprule
\textbf{Scenario} & \textbf{RL's role} & \textbf{Key methods} \\
\midrule
From scratch & Full policy learning & DQN, PPO, SAC \\
LLM alignment & Preference fine-tuning & PPO+RM, GRPO, DPO (+LoRA) \\
Diffusion / Flow & Denoising-chain MDP & DDPO, DPPO, DRaFT \\
Agent RL & Multi-round closed-loop & PPO, GRPO for LLM agents \\
Sim $\leftrightarrow$ Real & Transfer loop & Domain rand., world models \\
\bottomrule
\end{tabular}
\end{table}

The common thread: RL is \emph{preference optimisation under biased sampling}.
Whether the policy is a Q-network, an LLM, or a diffusion model, the same closed-loop challenge applies---the algorithm changes its own data distribution.
Every technique in this chapter is a different strategy for managing that fundamental difficulty in a different application context.
